%% Font size %%
\documentclass[11pt]{article}

%% Load the custom package
\usepackage{Mathdoc}

%% Numéro de séquence %% Titre de la séquence %%
\renewcommand{\centerhead}{S01 - Fractions - Encadrement entre deux entiers consécutifs}

%% Spacing commands %%
\renewcommand{\baselinestretch}{1} \setlength{\parindent}{0pt}

\begin{document}

\begin{exercice}[1][Exercice corrigé]
On écrit la fraction comme la somme d'un entier et d'une fraction plus petite que $1$ : c'est sa \textbf{décomposition} (exemple : $\dfrac{23}{6} = 3 + \dfrac{5}{6}$). La fraction est alors comprise entre cet entier et son \textbf{successeur}.
\vspace{.4em}

\textbf{Exemple :} $\dfrac{17}{5} = 3 + \dfrac{2}{5}$ (car $5 \times 3 = 15$ et $5 \times 4 = 20$), donc $3 \leqslant \dfrac{17}{5} \leqslant 4$.

\textbf{À toi :} donne la décomposition de chaque fraction, puis son encadrement.
\begin{multicols}{2}
\begin{enumerate}[itemsep=.6em]
\item $\dfrac{23}{4} = \ldots + \dfrac{\ldots}{4}$, donc $\ldots \leqslant \dfrac{23}{4} \leqslant \ldots$
\item $\dfrac{34}{6} = \ldots + \dfrac{\ldots}{6}$, donc $\ldots \leqslant \dfrac{34}{6} \leqslant \ldots$
\item $\dfrac{47}{9} = \ldots + \dfrac{\ldots}{9}$, donc $\ldots \leqslant \dfrac{47}{9} \leqslant \ldots$
\item $\dfrac{29}{8} = \ldots + \dfrac{\ldots}{8}$, donc $\ldots \leqslant \dfrac{29}{8} \leqslant \ldots$
\end{enumerate}
\end{multicols}
\end{exercice}

\begin{exercice}[1][Dénominateurs plus petits que 10]
Encadre chaque fraction entre deux entiers consécutifs.
\begin{multicols}{2}
\begin{enumerate}[itemsep=.4em]
\item $\ldots \leqslant \dfrac{13}{3} \leqslant \ldots$
\item $\ldots \leqslant \dfrac{19}{5} \leqslant \ldots$
\item $\ldots \leqslant \dfrac{26}{7} \leqslant \ldots$
\item $\ldots \leqslant \dfrac{11}{2} \leqslant \ldots$
\item $\ldots \leqslant \dfrac{37}{6} \leqslant \ldots$
\item $\ldots \leqslant \dfrac{22}{9} \leqslant \ldots$
\end{enumerate}
\end{multicols}
\end{exercice}

\begin{exercice}[1][Dénominateurs plus petits que 10]
Encadre chaque fraction entre deux entiers consécutifs.
\begin{multicols}{2}
\begin{enumerate}[itemsep=.4em]
\item $\ldots \leqslant \dfrac{17}{4} \leqslant \ldots$
\item $\ldots \leqslant \dfrac{25}{6} \leqslant \ldots$
\item $\ldots \leqslant \dfrac{31}{8} \leqslant \ldots$
\item $\ldots \leqslant \dfrac{16}{3} \leqslant \ldots$
\item $\ldots \leqslant \dfrac{29}{9} \leqslant \ldots$
\item $\ldots \leqslant \dfrac{23}{5} \leqslant \ldots$
\end{enumerate}
\end{multicols}
\end{exercice}

\begin{exercice}[2][Dénominateurs un peu plus grands]
Encadre chaque fraction entre deux entiers consécutifs.
\begin{multicols}{2}
\begin{enumerate}[itemsep=.4em]
\item $\ldots \leqslant \dfrac{47}{11} \leqslant \ldots$
\item $\ldots \leqslant \dfrac{58}{12} \leqslant \ldots$
\item $\ldots \leqslant \dfrac{67}{13} \leqslant \ldots$
\item $\ldots \leqslant \dfrac{89}{14} \leqslant \ldots$
\item $\ldots \leqslant \dfrac{47}{15} \leqslant \ldots$
\item $\ldots \leqslant \dfrac{100}{11} \leqslant \ldots$
\end{enumerate}
\end{multicols}
\end{exercice}

\end{document}
